
（六）怎样正确使用标点符号

大家知道，标点符号是书面语言里的一种极为重要的辅助手段。一句话，一篇文章，如果标点符号用得不对，就会叫人看不明白，甚至使人误解其意。例如在有的中学生的作文里，有时碰上这样一类的句子：
(1) 我赞成他也赞成你怎么样？
(2) 团委副书记、政治老师张涛、年级组组长、班主任黄萍都先后到我家里来看我。
很明显，第一句缺少两个逗号，使人看了之后不明白句子的意思是“我赞成，他也赞成，你怎么样？”还是“我赞成他，也赞成你，怎么样？”第二句的原意是说张涛和黄萍两位老师都来看我，但是由于在一人担任的两种职务之间打上了顿号，就成了四个成分并列，使人误解有四个人来看我了，所以，要么去掉第一个顿号和第三个顿号，要么将第二个顿号改为逗号，同时在“黄萍”后面加上一个逗号，句子的意思才能明确起来。
看来，不论少用、多用或者错用了标点符号，都会使书面语言不能准确地、明白地起到“交流思想”的作用。这就需要我们在使用标点符号时认真推敲一番才行。

——125——
